\bookchapter{第三十一章：ショート動画撮影とコンテンツ制作の入門}

\chapterintro{2026年、短動画はもう“若者の玩具”ではなく、すべての人が情報を得て自己表現する基礎手段になった。生活を記録しようが、知識を共有しようが、副業をしようが、スマホが最高の創作ツールだ——カメラを買う必要もPremiereを学ぶ必要もない。一台のスマホで始められる。この章はゼロからのスタートを助ける。}

\booksection{一、撮る前に明確に：誰に見せるのか？}

多くの人が動画撮影で最初に犯す間違いは“スマホを手に取ってすぐ撮る”ことで、撮り終わっても誰に見せるか分からない。撮る前にまず三つの質問に答えよう：

\begin{enumerate}
  \item \textbf{視聴者は誰か？}——子供向け、同年代向け、お年寄り向けでは、内容のスタイルが全く異なる。
  \item \textbf{どんな価値を提供できるか？}——笑わせる（エンタメ）、学ばせる（知識）、食欲をそそる（グルメ）、温める（感情）。少なくとも一つ。
  \item \textbf{どのプラットフォームで継続更新できるか？}——Douyin/Kuaishou（総合エンタメ）、REDnote（ライフスタイル）、Bilibili（深い内容）、動画アカウント（ソーシャル圏）。一つのメインプラットフォームを選んで深耕し、最初から全プラットフォーム配信はしない。
\end{enumerate}

\booksection{二、スマホ撮影の基本設定——撮影前に必須の5件事}

\begin{enumerate}
  \item \textbf{解像度は1080P、フレームレートは30fpsに設定}：これは短動画プラットフォームの標準だ。4Kはファイルが大きすぎてアップロードが遅く、60fpsはSNSではメリットがなくファイルがより大きい。
  \item \textbf{レンズをきれいに拭く}：最初で最も重要。撮る前に服でレンズを拭くだけで、画質が一ランク上がる。
  \item \textbf{フォーカスと露出をロック}：ファインダー内で被写体を2秒長押しすると"AE/AFロック"の表示が出る。これで撮影中に移動しても画面が明るすぎたり暗すぎたりしない。
  \item \textbf{横撮りか縦撮りか？}——Douyin/Kuaishou/動画アカウント/REDnoteに投稿\textbf{縦撮り}（9:16）、Bilibili/YouTubeに投稿\textbf{横撮り}（16:9）。プラットフォームを考えてから方向を決める。
  \item \textbf{グリッド線をオンにする}：設定 > カメラ > グリッド。三分割法の構図の基礎で、画面の重心を素早く見つけるのに役立つ。
\end{enumerate}

\booksection{三、カメラワーク五手——画面を“動かす”}

初心者が動画撮影で最大の問題は画面が“静”すぎること——ずっと立ったまま動かずに撮る。以下の五つの基本カメラワークを覚えれば、画面にすぐプロ感が出る：

\begin{itemize}
  \item \textbf{プッシュ（Push）}：スマホをゆっくり被写体に近づける。"近づく"緊張感や親密感を生む。
  \item \textbf{プル（Pull）}：スマホをゆっくり被写体から遠ざける。ディテールから全景に引き、環境を示す。
  \item \textbf{パン（Pan）}：スマホの位置はそのままで左右に回転。風景や空間の展示に適している。
  \item \textbf{トラック（Track）}：スマホを平行移動させ、被写体についていく。手ぶれ補正が必要で、ジンバルを使うのがベスト。
  \item \textbf{フォロー（Follow）}：スマホが被写体の後ろまたは前について歩く。Vlogで最もよく使うカメラワーク。
\end{itemize}

\textbf{ブレないことが第一}：手ブレした映像は素人っぽく見える。両手持ち片手持ちより安定；歩く時は膝を少し曲げ“忍び足”で上下の揺れを減らす；予算があればスマホジンバル（DJI OMシリーズ約500-800元）を買うと、効果が劇的に変わる。

\booksection{四、編集入門——CapCut（剪映）で十分}

剪映（CapCutの中国版）は現在初心者に最も適したスマホ編集ツールで、無料、機能充実、テンプレート豊富。以下は基本ワークフロー：

\begin{enumerate}
  \item \textbf{素材をインポート}：剪映を開く > 作成開始 > 撮影済み動画クリップを選択してインポート。
  \item \textbf{カット編集}：タイムラインをドラッグし、各クリップの冒頭と末尾の不要部分を切り取る。原則：各セグメントは情報量のある部分だけ残し、長さより短さを優先。
  \item \textbf{トランジションを追加}：二つの動画クリップの間で"トランジション"ボタンをタップし、シンプルな効果を選ぶ。初心者は"ディゾルブ"か"フェード"を推奨、派手な3Dトランジションは避ける。
  \item \textbf{字幕を追加}：剪映には"字幕認識"機能があり、ワンタップで自動生成。精度は約90\%、手動で一度校正すればよい。
  \item \textbf{BGMを付ける}：剪映の音楽ライブラリに著作権フリーの音楽が大量にある。内容に合うリズムのBGMを選び、音量を動画元音声の1/3程度に調整。
  \item \textbf{エクスポート}：1080Pでエクスポート、フレームレート30fps、ビットレートは"推奨"を選べばよい。
\end{enumerate}

\booksection{五、サムネイルの重要性——70\%の人がサムネイルだけで見るか決める}

短動画のサムネイルは“看板”で、タップするかどうかを決める。サムネイル作成のポイント：

\begin{itemize}
  \item \textbf{サムネイルに一言入れる}：大きな文字で動画のコア内容を書く。例えば"3分でスマホ写真編集をマスター"。文字は大きく、色は明るく、背景とのコントラストをつける。
  \item \textbf{サムネイルと内容を一致させる}："タイトル詐欺"をしない——サムネイルで一つのことを言っておきながら中身が違うと、視聴者はすぐスクロールし、アルゴリズムの評価が下がる。
  \item \textbf{人物の表情を生き生きと}：サムネイルに人がいる場合、生き生きとしたまたは誇張された表情の瞬間を選ぶ。人の顔は文字よりクリックを惹きつける。
  \item \textbf{剪映でサムネイルを作る}：剪映で最も良い一フレームを選び、テキストを重ねて、画像としてエクスポートしサムネイルにする。
\end{itemize}

\booksection{六、コンテンツ創作の三つの心法}

\begin{itemize}
  \item \textbf{最初の3秒で生死が決まる}：短動画プラットフォームでは、視聴者は3秒以内にスクロールするか見続けるかを決める。最初の3秒で"フック"を出す——一つの問い、一つのサスペンス、一つの意外な映像。
  \item \textbf{一本の動画は一つのことだけ話す}：初心者が最も犯しやすい間違いは一本の動画に内容を詰め込みすぎること。一本の動画は一つの問題だけを解決し、網羅より深掘りが大事。
  \item \textbf{才能より継続が大事}：アルゴリズム推薦には"アカウント育成"が必要。最初の20本の動画はほぼ再生数がないが、これは正常だ。毎週2-3本更新し続け、アルゴリズムに内容の方向性を認識させれば、トラフィックは徐々に伸びる。
\end{itemize}

\begin{tcbpractice}{実践：30秒の短動画を一本撮る}
\begin{enumerate}
  \item 得意な小ネタを一つ選ぶ（例："スマホの文字を大きくする方法"）。
  \item スマホ縦画面で3-5クリップ撮影、各5-10秒。
  \item 剪映で結合、カット、字幕追加、BGM付与。
  \item 大きな文字入りサムネイル画像を作る。
  \item エクスポート後、まず友人に見せてフィードバックを集めてから公開する。
\end{enumerate}

\end{tcbpractice}

\begin{tcbsummary}{本章のまとめ}
\begin{itemize}
  \item 撮る前に明確に：視聴者は誰か、どんな価値を提供するか、どのプラットフォームに投稿するか
  \item 撮影前の5件事：1080P/30fps、レンズ拭き、フォーカス露出ロック、横縦の確認、グリッドオン
  \item カメラワーク五手：プッシュ、プル、パン、トラック、フォロー——ブレないことが第一
  \item 編集は剪映：カットトランジション字幕BGMエクスポート、六ステップ
  \item 最初の3秒で生死が決まる、一本の動画は一つのことだけ、才能より継続が大事
\end{itemize}

\end{tcbsummary}
